8. Considere el elipsoide \( E=\left\{(x, y, z) \in \mathbb{R}^{3} \left\lvert\, \frac{x^{2}}{a^{2}}+\frac{y^{2}}{b^{2}}+\frac{z^{2}}{c^{2}} \leq 1\right.\right\} \), con \( a, b \) y \( c \) números reales positivos. \\
a) Calcule el volúmen de \( E \). \\
\aw Usamos el cambio de coordenadas esféricas con una sutileza en el coeficiente de cada variable.
\begin{align*}
    x &= a\rho sin\phi cos\theta \\
y &= b\rho sin\phi sin\theta \\
z &= c\rho cos\phi
\end{align*}

Donde $\rho, \phi \text{  y  } \theta$ son las coordenadas esféricas. Dado que trabajamos con una elipsoide nos gustaría trabajar con un problema que ya conocemos, las coordenadas esféricas, para ello necesitamos que sea una "esfera" que crece a diferentes velocidades para radios de cero a uno en las distintas 3 coordenadas, definimos el radio $\rho$ que llegue hasta uno pero dadas las condiciones en cada coordenada, en x se extenderá hasta un "radio a" (entre comillas porque la noción de radio está degenerada, es más bien un eje), en y radio b y en z radio c; por lo que nuestros parámetros se limitan como sigue:
\begin{itemize}
    \item $\rho$ varía de 0 a 1
    \item $\phi$ varía de 0 a $\pi$
    \item $\theta$ varía de 0 a 2$\pi$
\end{itemize}

El jacobiano de esta transformación es muy similar al de coordenadas esféricas "ordinarias", por propiedades de los determinantes podemos "factorizar" el término de una fila extrayéndolo del determinante. Al hacer esto tenemos el jacobiano de la transformación esférica que ya conocemos:
\[
J = \begin{vmatrix}
a \sin \phi \cos \theta & a \rho \cos \phi \cos \theta & -a \rho \sin \phi \sin \theta \\
b \sin \phi \sin \theta & b \rho \cos \phi \sin \theta & b \rho \sin \phi \cos \theta \\
c \cos \phi & -c \rho \sin \phi & 0
\end{vmatrix} = 
abc\begin{vmatrix}
\sin \phi \cos \theta & \rho \cos \phi \cos \theta & - \rho \sin \phi \sin \theta \\
\sin \phi \sin \theta & \rho \cos \phi \sin \theta & \rho \sin \phi \cos \theta \\
\cos \phi & -\rho \sin \phi & 0
\end{vmatrix} 
\]
\[
J =
abc \rho^2 \sin \phi
\]
Por lo tanto, el volumen del elipsoide es:
\[
V = \int_{0}^{2\pi} \int_{0}^{\pi} \int_{0}^{1} abc \rho^2 \sin \phi \, d\rho \, d\phi \, d\theta
\]
Resolvemos la integral separándola en 3, usando el mismo argumento de incisos anteriores pero generalizado para funciones de la forma $f(x,y,z)=f_1(x)f_2(y)f_3(z)$.\\

Integral respecto a $\rho$:
\[
  \int_{0}^{1} \rho^2 \, d\rho = \left[ \frac{\rho^3}{3} \right]_{0}^{1} = \frac{1}{3}
  \]


Integral respecto a $\phi$:
\[
  \int_{0}^{\pi} \sin \phi \, d\phi = \left[ -\cos \phi \right]_{0}^{\pi} = 2
  \]


Integral respecto a $\theta$:
\[
  \int_{0}^{2\pi} \, d\theta = 2\pi
  \]


Por lo tanto, el volumen es:
\[
V = abc \cdot \frac{1}{3} \cdot 2 \cdot 2\pi = \frac{4}{3} \pi abc
\]

b) Resuelva \( \iiint_{E}|x y z| d x d y d z \). \\
\aw Usamos la misma transformación:
\[
\iiint_{E} |x y z| \, dx \, dy \, dz = \int_{0}^{2\pi} \int_{0}^{\pi} \int_{0}^{1} |a b c \rho^3 \sin^2 \phi \cos \theta \sin \theta \cos \phi| \cdot abc \rho^2 \sin \phi \, d\rho \, d\phi \, d\theta
\]
Simplificamos la integral:
\[
= a^2 b^2 c^2 \int_{0}^{2\pi} \int_{0}^{\pi} \int_{0}^{1} \rho^5 \sin^3 \phi |\cos \theta \sin \theta \cos \phi| \, d\rho \, d\phi \, d\theta
\]

Separamos nuevamente; integral respecto a $\rho$:
\[
   \int_{0}^{1} \rho^5 \, d\rho = \left[ \frac{\rho^6}{6} \right]_{0}^{1} = \frac{1}{6}
   \]


Integral respecto a $\phi$:
\[
   \int_{0}^{\pi} \sin^3 \phi |\cos \phi| \, d\phi
   \]

Dividimos en dos partes, de \(0\) a \(\frac{\pi}{2}\) y de \(\frac{\pi}{2}\) a \(\pi\), ya que \(|\cos \phi|\) cambia de signo.

\[
   = 2 \int_{0}^{\frac{\pi}{2}} \sin^3 \phi \cos \phi \, d\phi = 2 \cdot \frac{1}{4} = \frac{1}{2}
   \]


Integral respecto a $\theta$:
\[
   \int_{0}^{2\pi} |\cos \theta \sin \theta| \, d\theta = \int_{0}^{2\pi} \frac{1}{2} |\sin(2\theta)| \, d\theta
   \]

Evaluamos desde 0 a $\pi$, multiplicamos por 2 por simetría:
\[
   = 2 \cdot \int_{0}^{\pi} \frac{1}{2} \sin(2\theta) \, d\theta = 2 \cdot \left[-\frac{1}{4} \cos(2\theta)\right]_{0}^{\pi} = 2 \cdot \frac{1}{2} = 1
   \]



Por lo tanto, el resultado de la integral es:
\[
a^2 b^2 c^2 \cdot \frac{1}{6} \cdot \frac{1}{2} \cdot 1 = \frac{a^2 b^2 c^2}{12}
\]